\documentclass[11pt]{beamer}

\usetheme{default}
\usecolortheme{default}
\beamertemplatenavigationsymbolsempty
\setbeamertemplate{footline}[frame number]

\usepackage[T1]{fontenc}
\usepackage[utf8]{inputenc}
\usepackage{listings}
\usepackage{xcolor}
\usepackage{tikz}
\usepackage{tikz-cd}
\usepackage{mathtools}

\hypersetup{
    colorlinks = true
}

\AtBeginSection[]{
\begin{frame}{Contents}
    \tableofcontents[currentsection]
\end{frame}
}


\lstdefinestyle{shell}{
  basicstyle=\ttfamily\footnotesize,
%  identifierstyle=\ttfamily
  breaklines=true,
  showstringspaces=false,
  frame=single,
  framerule=0.4pt,
  xleftmargin=0.5em,
  xrightmargin=0.5em,
  columns=flexible,                   % Removes fixed-column spacing
  showstringspaces=false                  % Hides space markers in strings
}
\lstset{style=shell}

\title{Git and GitHub: an introduction}
\subtitle{Lecture 2 --- Branching, Merging, and GitHub}
\author{}
\date{}

\newcommand{\exercise}[1]{
  \begin{frame}{Exercise}
    #1
  \end{frame}
}

\begin{document}

\begin{frame}
  \titlepage
\end{frame}

\begin{frame}{Outline}
  \tableofcontents
\end{frame}

% ------------------------------------------------------------------
\section{Recap}

\begin{frame}[fragile]{Where we stopped}
  The basic Git cycle:
\begin{lstlisting}
git status
git add <file>
git commit
git log
\end{lstlisting}
  \vfill
  Today:
  \begin{itemize}
    \item Branches and merging
    \item GitHub accounts and SSH keys
    \item Pushing, pulling, collaborating
  \end{itemize}
\end{frame}

% ------------------------------------------------------------------
\section{Branches}

\begin{frame}[fragile]{What is a branch?}
  The history of a repository is a graph of commits.
  %
  So far, the history looked like this:
  \begin{flushleft}
    \begin{tikzcd}
      \cdots
      \arrow[r]
      &
      A
      \arrow[r]
      &B
      \arrow[r]
      &C
      \arrow[r]
      &
      \underset{\mathtt{master}}{D}
    \end{tikzcd}
  \end{flushleft}
  If more than one branch is present, the history may look like this:
  \begin{flushleft}
    \begin{tikzcd}
      \cdots
      \arrow[r]
      &
      A
      \arrow[r]
      \arrow[rd]
      &B
      \arrow[r]
      &C
      \arrow[r]
      &
      \underset{\mathtt{master}}{D}
      \\
      &
      &
      X
      \arrow[r]
      &
      \underset{\mathtt{slave}}{Y}
    \end{tikzcd}
  \end{flushleft}
  \vfill
  A \emph{branch} is just a movable pointer to one of these commits.

  % \vfill
  The default branch is usually called \texttt{master}.
\end{frame}

\begin{frame}{Why branches?}
  Typical uses:
  \begin{itemize}
    \item Try an alternative proof without touching the main version
    \item Experiment with a different structure of a paper
%    \item Develop a new feature in code while keeping a working version
    \item Let a coauthor work in parallel without stepping on each other
  \end{itemize}
  \vfill
  If the experiment fails: delete the branch, no harm done.
  \\
  If it works: merge it back.
  \vfill
  Takeaway: branches are cheap.  Use them!
\end{frame}

\begin{frame}[fragile]{Working with branches}
\begin{lstlisting}
git branch                     # list branches
git branch new-proof           # create a branch
git branch                     # see the new branch
git switch new-proof           # switch to it
git switch master              # go back
git diff master new-proof      # compare to branches
\end{lstlisting}
%git switch -c new-proof       # create and switch
%git branch -d new-proof       # delete (if merged)
%  \vfill
%  \small (Older Git versions use \texttt{git checkout} instead of \texttt{git switch}.)
\end{frame}

% ------------------------------------------------------------------
\section{Merging}

\begin{frame}[fragile]{Merging a branch}
  From the branch you want to merge \emph{into}:
\begin{lstlisting}
git switch master
git merge new-proof
\end{lstlisting}
  \vfill
  Two cases:
  \begin{itemize}
  \item
    \textbf{Fast-forward}: \texttt{master} had no new commits;
    Git just moves the pointer.
  \item
    \textbf{Three-way merge}: both branches advanced;
    Git creates a new ``merge commit''.
  \end{itemize}
\end{frame}

\begin{frame}[fragile]{Fast forward}
   \texttt{master} had no new commits;
   Git just moves the pointer.
   \vfill

   Before:
  \begin{flushleft}
    \begin{tikzcd}
      \cdots
      \arrow[r]
      &
      A
      \arrow[r]
      &B
      \arrow[r]
      &
      \underset{\mathtt{master}}{C}
      \arrow[rd]
      \\
      &
      &
      &&
      X
      \arrow[r]
      &
      \underset{\mathtt{slave}}{Y}
    \end{tikzcd}
  \end{flushleft}

  After:
  \begin{flushleft}
    \begin{tikzcd}
      \cdots
      \arrow[r]
      &
      A
      \arrow[r]
      &B
      \arrow[r]
      &
      C
      \arrow[rd]
      \\
      &
      &
      &&
      X
      \arrow[r]
      &
      \underset{\substack{\mathtt{master}\\ \mathtt{slave}}}{Y}
    \end{tikzcd}
  \end{flushleft}
\end{frame}

\begin{frame}[fragile]{Three-way merge}
  Both branches advanced;
    Git creates a new ``merge commit''.
    \vfill
    Before: 
  \begin{flushleft}
    \begin{tikzcd}
      \cdots
      \arrow[r]
      &
      A
      \arrow[r]
      \arrow[rd]
      &B
      \arrow[r]
      &C
      \arrow[r]
      &
      \underset{\mathtt{master}}{D}
      \\
      &
      &
      X
      \arrow[r]
      &
      \underset{\mathtt{slave}}{Y}
    \end{tikzcd}
  \end{flushleft}

  After:
  \begin{flushleft}
    \begin{tikzcd}
      \cdots \arrow[r] & A \arrow[r] \arrow[rd] &B \arrow[r] &C
      \arrow[r] & D \arrow[r] & \underset{\substack{\mathtt{master}\\
          \mathtt{slave}}}{U}
      \\
      & & X \arrow[r] & Y \arrow[rru]
    \end{tikzcd}
  \end{flushleft}
\end{frame}

\begin{frame}[fragile]{Merge conflicts}
  If both branches changed \emph{the same lines}, Git cannot decide.
  It marks the conflict in the file:
\begin{lstlisting}
<<<<<<< HEAD
The proof uses the dominated convergence theorem.
=======
The proof uses Fatou's lemma.
>>>>>>> new-proof
\end{lstlisting}
  \vfill
  To resolve:
  \begin{enumerate}
    \item Edit the file, keep what you want, remove the markers.
    \item \texttt{git add <file>}
    \item \texttt{git merge -{}-continue}
  \end{enumerate}
\end{frame}

\exercise{
  \begin{enumerate}
  \item
    \begin{itemize}
    \item Create a new branch and switch to it
    \item Make a new commit
    \item Switch back to \texttt{master} and merge the new branch
    \end{itemize}
  \item
    \begin{itemize}
    \item Create a new branch and switch to it
    \item Create a new file and commit the change
    \item Switch back to \texttt{master} and edit the old file
    \item Merge the new branch
    \end{itemize}
  \item
    \begin{itemize}
    \item Create a new branch and switch to it
    \item Make new commit
    \item Switch back to \texttt{master}
    \item Edit the file, in the same place, where it is modified in
      the other branch
    \item Merge the new branch
    \end{itemize}
  \end{enumerate}
}

\begin{frame}[fragile]{Note on \LaTeX}
  Git works at the level of \emph{lines} (in the source code).

  When writing in \LaTeX, keep each sentence in different
  lines.
  
  Better yet, keep lines short ($\leq 75$) and separate sentences
  with \texttt{\%}

  \vfill
  {\color{red} Bad:}
\begin{lstlisting}
A causal vector $v$ will be called future pointing if $g(v,X)>0$.  Past pointing otherwise.
\end{lstlisting}
  {\color{blue} Good:}
\begin{lstlisting}
A causal vector $v$ will be called future pointing if $g(v,X)>0$.
Past pointing otherwise.
\end{lstlisting}
  {\color{blue} Better:}
\begin{lstlisting}
A causal vector $v$ will be called future pointing
if $g(v,X)>0$.
%
Past pointing otherwise.
\end{lstlisting}
\end{frame}


\begin{frame}{A note on rebasing}
  You may hear about \texttt{git rebase}.
  \vfill
  It is an alternative way to integrate changes, with a cleaner history
  but sharper edges.
  \vfill
  \alert{We will not cover it.} Stick to \texttt{merge} for now;
  you can learn rebasing later.
\end{frame}

% ------------------------------------------------------------------
\section{GitHub}

\begin{frame}{What GitHub gives you}
  \begin{itemize}
    \item Online hosting for your repositories
    \item Backup (your laptop is not eternal)
    \item A way to share code and papers
    \item Collaboration tools: pull requests, issues\ldots
    \item Easy web interface
  \end{itemize}
  \vfill
  Free for public \emph{and} private repos.
%  \\
 % Students get extra perks via \emph{GitHub Education}.
\end{frame}

% \begin{frame}{Public vs.\ private}
%   \textbf{Public}
%   \begin{itemize}
%     \item Anyone can see the code
%     \item Good for published work, open-source projects
%   \end{itemize}
%   \vfill
%   \textbf{Private}
%   \begin{itemize}
%     \item Only you and invited collaborators see it
%     \item Good for work in progress, unpublished papers
%   \end{itemize}
%   \vfill
%   You can change visibility later.
% \end{frame}

% ------------------------------------------------------------------
\section{SSH keys}

\begin{frame}{Why SSH?}
  To push to GitHub you must prove who you are.
  \vfill
  Two options:
  \begin{itemize}
    \item HTTPS + personal access token (awkward)
    \item \textbf{SSH keys} (recommended)
  \end{itemize}
  \vfill
  An SSH key is a pair:
  \begin{itemize}
    \item a \emph{private} key --- stays on your computer, never shared
    \item a \emph{public} key --- uploaded to GitHub
    \end{itemize}
    \vfill
    Three steps to enable SSH login to GitHub:
    \begin{itemize}
    \item Generate your SSH key pair
    \item Upload your public SSH key
    \item Trust GitHub's public keys
    \end{itemize}
\end{frame}

\begin{frame}[fragile]{Generating a key pair}
\begin{lstlisting}
ssh-keygen
\end{lstlisting}
  Press Enter to accept the default location.
  Choose a passphrase (leave blank for no passphrase).
  \vfill
  Two files are created:
\begin{lstlisting}
~/.ssh/id_ed25519          # PRIVATE -- never share
~/.ssh/id_ed25519.pub      # PUBLIC  -- upload to GitHub
\end{lstlisting}
\end{frame}

\begin{frame}[fragile]{Uploading the public key}
  Copy the content of the \texttt{.pub} file:
\begin{lstlisting}
cat ~/.ssh/id_ed25519.pub
\end{lstlisting}
  On GitHub:
  \begin{enumerate}
    \item Settings $\to$ SSH and GPG keys $\to$ New SSH key
    \item Paste the content, give it a title (e.g., ``laptop'',
      ``work-pc'' \dots)
  \end{enumerate}
\end{frame}

\begin{frame}[fragile]{Trust GitHub's public keys}

  Try to connect to GitHub using SSH{}:
\begin{lstlisting}
ssh -T git@github.com
\end{lstlisting}
You will be asked if you trust the key; answer `yes'.

  \vfill

  You should see a message greeting you by username.

  \vfill

  You need to setup SSH keys on every computer you use.


\end{frame}

% ------------------------------------------------------------------
\section{Remotes}

\begin{frame}[fragile]{Connecting a local repo to GitHub}
  On GitHub: create an empty repository (no README, no license).
  \vfill
  Back in your terminal:
\begin{lstlisting}
git remote add origin git@github.com:you/repo.git
git remote -v                        # List remote repositories
\end{lstlisting}
  \texttt{origin} is just a nickname for the remote; you can have
  several.
  \vfill
  Next you need to connect a specific branch to a remote repository
\begin{lstlisting}
git switch <branch>
git push -u origin        # Tell git to push <branch> to origin
\end{lstlisting}
You need to do this for all branches
\end{frame}

\begin{frame}[fragile]{Push and pull}
\begin{lstlisting}
git pull              # fetch and merge remote commits
git push              # send local commits to GitHub
\end{lstlisting}
%git fetch             # fetch only, do not merge
  \vfill
  Rule of thumb:
  \begin{itemize}
  \item \texttt{git pull} before you start working
  \item \texttt{git push} when you are done
  \end{itemize}
  \vfill
  Daily workflow:
\begin{lstlisting}
git switch your-branch
git pull
# standard git edit cycle (edit; git add; git commit)
git push
\end{lstlisting}
\end{frame}

\exercise{
  \begin{itemize}
  \item Create an empty repository on GitHub
  \item On your computer add local the remote GitHub repository
  \item Push your local master branch
  \item Make a few commits on a different branch and push them
  \item Inspect the history on GitHub
  \end{itemize}
}


% ------------------------------------------------------------------
\section{Collaboration}

\begin{frame}{A typical collaborative workflow}
  \begin{enumerate}
    \item \texttt{git clone} the repository (or \texttt{git pull} if
      you are already collaborating)
    \item Create a branch for your change (do not touch other people's branches)
    \item Commit and push the branch
    \item Open a \emph{pull request} on GitHub
    \item Discuss, review, fix
    \item Merge the pull request
    \item Pull the updated \texttt{master} locally
  \end{enumerate}
  \vfill
  This scales from two people to thousands.
\end{frame}

\begin{frame}{Pull requests in short}
  A \emph{pull request} (PR) is a GitHub feature, not a Git command.
  \vfill
  It is a proposal: \emph{``please merge my branch into yours''}.
  \vfill
  It lets collaborators:
  \begin{itemize}
    \item see a diff of the proposed change
    \item comment line by line
    \item request modifications
  \end{itemize}
\end{frame}

\exercise{
  To be done by user $A$ and $B$
  \begin{itemize}
  \item[$A$] Give access to $B$ to your private repository
  \item[$B$] Clone the repository of $A$
  \item[$B$] Create a new branch and make some edits
  \item[$B$] Push the new branch
  \item[$B$] Open a pull request
  \item[$A$] Review the pull request, discuss with $B$, and merge it
  \end{itemize}

  Repeat the exercise swapping the roles
}

% ------------------------------------------------------------------
\section{Hands-on}

\begin{frame}{Exercise}
  In pairs:
  \begin{enumerate}
    \item Each of you creates a repo on GitHub and pushes the project
          from Lecture 1.
    \item Clone your partner's repo.
    \item Create a branch, add a small change (e.g.\ a new paragraph).
    \item Push the branch and open a pull request.
    \item As the owner: review, comment, and merge the {PR}.
    \item Pull the updated \texttt{master} locally.
  \end{enumerate}
\end{frame}


\end{document}
